\section{Non-self-dual categories}\label{sec:ns}

\subsection{Statement and proof}

The non-self-dual sphericality proof uses no classification theorem,
fusion-ring census, or multiplicity bound.


\begin{theorem}\label{ns:thm:ns-spherical}
Every complex fusion category of rank four with a non-self-dual simple
object admits a spherical structure.
\end{theorem}

We prove the theorem using the sphericalization and formal codegrees.
The construction and its section criterion are recalled in
\cite[Section~5.1]{Ostrik3}; see also \cite[Section~5.3]{Bartlett}.
Write the simple objects as $1,X,Y,Z$, where $X^*=Z$ and $Y^*=Y$.
Let $R=K_0(\C)$, let $\mathcal S$ be the sphericalization, and
let $\varepsilon$ be its central sign object. Its two quotients are
\[
 R=K_0(\mathcal S)/(\varepsilon-1),\qquad
 T=K_0(\mathcal S)/(\varepsilon+1).
\]
Choose dual-compatible simple lifts. If $N_{ij}^k$ are the coefficients
of $R$ and $t_{ij}^k$ those of $T$, then
\begin{equation}\label{ns:eq:ns-domination}
 |t_{ij}^k|\leq N_{ij}^k,\qquad t_{ij}^k\equiv N_{ij}^k\pmod2.
\end{equation}
These are respectively the sum and difference of the two lift
multiplicities. Changing a lift changes its basis sign in $T$.

\begin{lemma}\label{ns:lem:ns-fields}
Suppose that $\C$ admits no spherical structure. There is a real quadratic
field $K$ and imaginary quadratic fields $L,M$ such that
\[
 R\otimes\Q\cong K\times L,\qquad
 T\otimes\Q\cong K\times M.
\]
The formal codegrees are, respectively,
\[
 D,D',f,f\qquad\text{and}\qquad E,E',g,g,
\]
where $D=\operatorname{FPdim}(\C)$ is irrational,
$f,g\in\ZZ_{\geq3}$, and primes denote the nontrivial embedding
of $K$. We can choose the lifts and a real character $\beta$ of $T$
so that all its values are positive, its codegree is $E$, and
\begin{equation}\label{ns:eq:ns-basic}
 E\geq E',\qquad E<D,\qquad
 E/D\ \text{and}\ E/g\ \text{are algebraic integers}.
\end{equation}
Every displayed codegree is greater than one, and $f\leq D$.
\end{lemma}

\begin{proof}
Lemma~\ref{core:lem:rank-four-algebra} applies to both regular
based $*$-algebras, so both quotients and their direct sum
$K_0(\mathcal S)\otimes\CC$ are commutative.

The complex-linear based duality has trace two, since it fixes $1,Y$
and exchanges $X,Z$. On primitive idempotents it permutes the
characters. Exactly two characters are therefore fixed by duality.
The positive Hermitian pairing gives
$\chi(b^*)=\overline{\chi(b)}$, so these are precisely the two real
characters. In particular, the Galois orbit of any real character
has size at most two: the equations $\chi(b)=\chi(b^*)$ persist under
Galois automorphisms.

If $D$ were rational, $\C$ would be weakly integral and hence
pseudo-unitary and spherical by \cite[Proposition~8.24]{ENO}.
Otherwise the Frobenius--Perron character field is quadratic and
equals $K=\Q(D)$. The dimension-field theorem
\cite[Corollary~2.15]{Ostrik3}, applied to $\mathcal S$, puts $K$
inside its spherical dimension field. That field has degree at most
two by the preceding character argument, so it too is $K$.
The remaining two-dimensional factors have two nonreal characters
and are imaginary quadratic fields. Their conjugate characters have
the same real formal codegree, hence a rational algebraic integer;
these are $f$ and $g$.

For a commutative algebra with the above orthonormal based pairing,
the formal codegrees are positive and satisfy
$\sum_\chi f_\chi^{-1}=1$. Positivity follows also from
$f_\chi=\sum_i|\chi(b_i)|^2$. As there are four characters, each
codegree is strictly greater than one. The repeated codegrees
therefore satisfy $f,g>2$. Character spectral domination gives
$|\chi(b_i)|\leq\operatorname{FPdim}(b_i)$ for $R$, whence $f\leq D$.

Initially write $\delta$ for the codegree of the spherical dimension
character of $T$, so $\dim(\mathcal S)=2\delta$. The formal codegrees
of $\mathcal S$ are twice those of $R$ and $T$. The induction theorem
\cite[Theorem~2.13]{Ostrik3} implies that
$\delta/D$, $\delta/D'$ and $\delta/g$ are algebraic integers.
Conjugate only the signed dimension character, if necessary, to
choose its codegree $E\geq E'$. The three divisibility statements
show that $E/D$ and $E/g$ remain integral after either embedding.
Every value of the conjugated character is real and nonzero. Choose
lift signs to make these values positive; dual compatibility is
preserved because the character agrees on dual basis elements.
This does not change the category or its sphericalization.

By \eqref{ns:eq:ns-domination}, spectral domination bounds each positive
$\beta(b_i)$ by $d_i=\operatorname{FPdim}(b_i)$. Thus
$E=\sum_i\beta(b_i)^2\leq\sum_i d_i^2=D$. If equality holds, every
$\beta(b_i)=d_i$, and the two dimension equations give
\[
 \sum_k (N_{ij}^k-t_{ij}^k)d_k=0.
\]
All summands are nonnegative, so $t_{ij}^k=N_{ij}^k$ for every triple.
The unselected lift consequently has multiplicity zero in every
product of selected lifts. They form a fusion subcategory of
$\mathcal S$, on which the forgetful functor is an equivalence with
the given category $\C$. Restriction transports a spherical
structure to $\C$, a contradiction. Hence $E<D$.
\end{proof}

\begin{lemma}\label{ns:lem:ns-small-codegree}
In the situation of Lemma~\ref{ns:lem:ns-fields}, put
\[
 n=DD',\qquad k=\frac{EE'}{DD'}.
\]
Then $k\in\{1,2\}$ and $g\in\{3,4\}$, with $k=2$ possible only
when $g=3$.
\end{lemma}

\begin{proof}
The norm of the totally positive algebraic integer $E/D$ makes $k$
a positive integer. Set
\[
 p=1-\frac2f=\frac1D+\frac1{D'},\qquad
 q=1-\frac2g=\frac1E+\frac1{E'}.
\]
Since $E\geq\sqrt{kn}$ and $E<D$, strict monotonicity of
$t+kn/t$ above its minimum gives
\begin{equation}\label{ns:eq:ns-key}
 kq<p+\frac{k-1}{D}\leq1+\frac{k-3}{D}.
\end{equation}
Here we used $f\leq D$. If $k\geq3$, then $q\geq1/3$ and $D\geq4$
give the contradiction
\[
 \frac{k}{3}<1+\frac{k-3}{D}\leq\frac{k+1}{4}\leq\frac{k}{3}.
\]
Thus $k\leq2$. If $k=2$ and $g\geq4$, the left side of
\eqref{ns:eq:ns-key} is at least one, whereas its right side is
$1-1/D<1$.

If $k=1$, put $r=D/E$. It is the inverse of an algebraic-integer
unit of norm one, hence is itself integral. It is totally positive,
and $r>1$, so $r+r^{-1}$ is an integer at least three.
Furthermore $E'=rD'$ and $E\geq E'$ imply $D\geq r^2D'$. Therefore
\[
 \frac qp=\frac{D/r+rD'}{D+D'}
 \leq\frac2{r+r^{-1}}\leq\frac23.
\]
The first inequality is equivalent to
$(r^2-1)(r^2D'-D)\leq0$. Since $p<1$, we get $q<2/3$ and $g<6$.

A nonreal character of $T$ has values
$(1,\alpha,t,\overline\alpha)$, with $t\in\ZZ$ and
$\alpha$ an imaginary quadratic algebraic integer. Put
$a=\operatorname{Tr}(\alpha)$ and $b=\operatorname{N}(\alpha)$.
Orthogonality to its conjugate character gives
\begin{equation}\label{ns:eq:ns-nonreal}
 2b=1+a^2+t^2,\qquad g=1+t^2+2b=2+a^2+2t^2.
\end{equation}
For $g=5$, the latter identity forces $a^2=t^2=1$, making $2b=3$,
impossible. Hence $g$ is three or four.
\end{proof}

\begin{proof}[Proof of Theorem~\ref{ns:thm:ns-spherical}]
Suppose that $\C$ admits no spherical structure, and use the preceding
lemmas. Unit, duality and Frobenius reciprocity give the signed table
\begin{equation}\label{ns:eq:ns-signed-table}
\begin{aligned}
 X^2&=P X+L Y+C Z,& XY&=Q X+J Y+L Z,\\
 Y^2&=1+J X+H Y+J Z,& YZ&=L X+J Y+Q Z,\\
 Z^2&=C X+L Y+P Z,& XZ&=1+P X+Q Y+P Z.
\end{aligned}
\end{equation}
All six parameters are integers. This is the signed version of the
standard non-self-dual table in \cite[Proposition~2.1]{Larson}.

If $g=3$, \eqref{ns:eq:ns-nonreal} gives $t=0$, $a=\pm1$, $b=1$.
Orthogonality to the positive character $\beta$ gives
$1+a\beta(X)=0$, hence $a=-1$, $\beta(X)=\beta(Z)=1$, and
$\alpha$ is a primitive cube root of unity. Evaluating
$X^2$, $XY$ and $Y^2$ in this character gives
\[
 C=J=1,\qquad P=L=Q=0,\qquad H=\mu\in\ZZ.
\]
Writing $v=\beta(Y)>0$, we obtain
\begin{equation}\label{ns:eq:ns-ng-dimensions}
 v^2=3+\mu v,\quad v'=-3/v,\quad
 E=3+v^2,\quad E'=3+9/v^2.
\end{equation}
The ordering $E\geq E'$ implies $v\geq\sqrt3$, so $\mu\geq0$.
In particular,
\begin{equation}\label{ns:eq:ns-ng-norm}
 E+E'=\mu^2+12,\qquad EE'=3(\mu^2+12).
\end{equation}

If $g=4$, \eqref{ns:eq:ns-nonreal} gives $a=0$, $t=\pm1$, $b=1$.
Orthogonality to $\beta$ gives $t=-1$ and $\beta(Y)=1$, with
$\alpha=\pm i$. Evaluation of \eqref{ns:eq:ns-signed-table} gives
\[
 L=1,\qquad J=Q=H=0,\qquad P=C=\mu\in\ZZ.
\]
Writing $u=\beta(X)=\beta(Z)>0$, we have
\[
 u^2=1+2\mu u,\quad u'=-1/u,\quad
 E=2+2u^2,\quad E'=2+2/u^2.
\]
Again $E\geq E'$ gives $\mu\geq0$. The case $\mu=0$ would make
$\beta$ rational, contrary to its field being $K$, so $\mu\geq1$.
Thus $u=\mu+\sqrt{\mu^2+1}>2$ and $E'<5/2$. Here $k=1$, so the
totally positive unit $r=D/E>1$ has integral trace at least three:
\[
 r\geq\frac{3+\sqrt5}{2}>\frac52>E'.
\]
This gives $D'=E'/r<1$, contrary to Lemma~\ref{ns:lem:ns-fields}.

It remains to exclude $g=3$. First suppose $\mu>0$. Integrality of
$E/3$ and \eqref{ns:eq:ns-ng-norm} give
$(\mu^2+12)/3\in\ZZ$, hence $3\mid\mu$ and $\mu\geq3$.
Moreover
\[
 K=\Q\bigl(\sqrt{\mu^2+12}\bigr),\qquad
 \mu^2+12=3\bigl(3(\mu/3)^2+4\bigr).
\]
The radicand has $3$-adic valuation one, so $K$ is neither
$\Q(\sqrt5)$ nor $\Q(\sqrt2)$.
From \eqref{ns:eq:ns-ng-dimensions}, $v>7/2$ and
$v^2=3+\mu v>27/2$, whence
\begin{equation}\label{ns:eq:ns-small-conjugate}
 E'<\frac{11}{3}<2+\sqrt3.
\end{equation}
For $k=1$, the unit $r=D/E>1$ cannot have trace three, since that
would generate $\Q(\sqrt5)$. Its trace is at least four,
so $r\geq2+\sqrt3>E'$, again giving $D'<1$.

For $k=2$, put $a_0=E/D$ and $s=2D/E$. The known algebraic integer
$a_0$ has norm two, so $s=2/a_0=a_0'$ is an algebraic integer of
norm two. It is totally positive and $s>2$. Its integral trace
$s+2/s$ is therefore greater than three. Trace four would generate
$\Q(\sqrt2)$, which is excluded; thus the trace is at least
five and $s>4>E'$. Since $D'=E'/s$, this too contradicts $D'>1$.

Finally suppose $\mu=0$. Then $E=E'=6$ and $n=DD'$ is either
$36$ or $18$. Denote the unsigned parameters corresponding to
$(C,H,J,L,P,Q)$ by $(c,e,k_0,\ell,p_0,q_0)$. Their parities are
\[
 c,k_0\ \text{odd},\qquad e,\ell,p_0,q_0\ \text{even}.
\]
For a nonreal character $(1,\alpha,t,\overline\alpha)$ of $R$,
the equations for $Y(X-Z)$ and $X^2-Z^2$ give
$t=q_0-\ell$ and $\alpha+\overline\alpha=p_0-c$.
Applying \eqref{ns:eq:ns-nonreal} to $R$ therefore gives
\[
 f=2+(p_0-c)^2+2(q_0-\ell)^2\equiv3\pmod8.
\]
But $\operatorname{Tr}(D)=n(1-2/f)$ is an integer, so $f\mid2n$.
Since $f$ is odd and $n\in\{18,36\}$, it divides nine; the
congruence and $f\geq3$ force $f=3$. For $n=36$ the trace of $D$
is $12$, giving $D=D'=6$, contrary to irrationality. For $n=18$
the trace is $6$, giving negative discriminant $36-72$,
contrary to reality. Every possibility has now been excluded.
\end{proof}



\begin{corollary}\label{ns:cor:pseudo-unitary-conjugate}
Every non-self-dual rank-four complex fusion category has a
pseudo-unitary Galois conjugate.
\end{corollary}
\begin{proof}
It is spherical by Theorem~\ref{ns:thm:ns-spherical}. Duality has
trace two on the based algebra and permutes primitive idempotents,
so exactly two characters are real. This condition is Galois stable.
If $D$ is rational, weak integrality gives pseudo-unitarity. Otherwise
the Frobenius--Perron orbit consists of both real characters: its
degree is two. The spherical dimension character is real and hence
belongs to this orbit. Apply Lemma~\ref{core:lem:galois-transport}.
\end{proof}

